% =============================================================================
% Snippet: Numerical Experiments and Rig Validation
% Target: the topological-bounds paper (references \ref{sec:topological_bounds})
%
% Every figure below is traceable to design/specs/calibration-rig-spec.md or to
% design/methodology/evidence/evidence.json in the ranking-residuals repo, where
% each is a registered claim with a tolerance. Where a value
% depends on a configuration default, the default is named inline, because two
% numbers in the first draft of this section had silently gone stale when a
% default moved.
% =============================================================================

\section{Numerical Experiments and Rig Validation}
\label{sec:numerical_experiments}

To validate the theoretical invariance properties established in
Section~\ref{sec:topological_bounds}, we evaluate the framework using a synthetic
Hodge calibration instrument ($1{,}586$ lines of rig code, exercised by a further
$571$ lines of acceptance tests). The instrument generates paired comparisons
across mixed populations consisting of a strictly ordered integer topology and an
unorderable, rotationally structured complex topology on the unit circle. The
decomposition itself is not reimplemented: the rig is a data source for a fixed
reference implementation it imports and never forks.

\subsection{Asymptotic Intercept Estimation and Logit Bias}
In finite-sample comparison regimes where each item pair is evaluated over a
fixed budget of $k$ comparisons, the observed edge flow $Y$ is contaminated by
low-$k$ logit bias. To decouple this sampling noise from the structural
unrankability floor, we model the measured harmonic energy $\Vert P_h Y \Vert^2$
using an ordinary least squares (OLS) linear regression in the inverse comparison
domain $1/k$:

\begin{equation}
\Vert P_h Y \Vert^2 = \text{floor} + \frac{c}{k}
  + \mathcal{O}\!\left(\frac{1}{k^2}\right)
\label{eq:intercept}
\end{equation}

The target invariant structural parameter is recovered exactly as the asymptotic
intercept where $1/k \to 0$. The slope $c$ is a fitted coefficient; it has an
independent closed-form prediction $c_{\text{oracle}}$ (Section~\ref{subsec:filling}),
and holding the two notationally distinct is what makes the agreement between
them usable as a misspecification check.

As illustrated in empirical sweeps, the small-$k$ tail is heavily dominated by
the non-linear $\mathcal{O}(1/k^2)$ logit-bias term. Fitting across the full
unwindowed grid ($k \in [8, 4096]$) artificially shifts the intercept upward to
$0.186$ against a true misspecification oracle of $\epsilon^2 = 0.090$; on the
shorter $k \in [8, 1024]$ grid the same fit gives $0.223$.
Restricting the estimation window to the clean asymptotic tail ($k \ge 64$)
suppresses the quadratic bias and recovers the analytical oracle floor to
$0.87\text{--}0.95\times$ across separation parameters $\beta \in [0.15, 0.30]$.
The decisive gain here is not the residual bias but its \emph{insensitivity} to
$\beta$: the unwindowed fit ranges over $0.83\text{--}2.48\times$ on the same
sweep, so the fixed cut converts a parameter-dependent estimate into a
parameter-robust one. Tightening the band further, to
$0.94\text{--}1.01\times$, requires the derived window of
Section~\ref{subsec:filling} rather than any fixed cut.

We note for completeness that modelling the contaminating term explicitly, by
admitting a $c_2/k^2$ regressor, is worse than discarding the points it
contaminates: the extended design is poorly conditioned and the additional
parameter absorbs the intercept ($0.58\text{--}0.90\times$).

\subsection{Topological Dependences of the Choice of Filling}
\label{subsec:filling}
The choice of topological simplicial closure significantly alters the internal
capacity of the harmonic space. We evaluate two distinct filling regimes across
identical boundary topologies:

\begin{enumerate}
    \item \textbf{Observed Filling:} Complex cyclic triads are closed into solid
    2D faces upon observation, yielding a tighter harmonic footprint
    ($b_1 = 2$, $c_{\text{oracle}} \approx 18$ on the reference graph).
    \item \textbf{Empty Filling:} Triads are left un-filled, treating local
    cyclic contradictions as persistent global harmonic holes
    ($b_1 = 20$, $c_{\text{oracle}} = 160$).
\end{enumerate}

Because $c_{\text{oracle}}$ scales by nearly an order of magnitude based on
filling choice, a fixed data-gathering threshold cannot guarantee numerical
stability: the $k \ge 64$ cut calibrated under the observed filling under-reads
the floor by a factor of six under the empty one ($0.0156$ against a true
$0.090$). We resolve this by implementing a dynamically derived minimum budget
window ($k_{\min}$), computed per configuration via

\begin{equation}
k_{\min} = \frac{c_{\text{oracle}}}{\rho \cdot \text{floor}},
\qquad
c_{\text{oracle}} = \operatorname{tr}\!\Big(P_h \cdot
  \operatorname{diag}\big(1/(p_e(1-p_e))\big)\Big),
\label{eq:kmin}
\end{equation}

where $\rho$ is a resolvability margin: at $k_{\min}$ the variance term equals
$\rho$ times the floor being resolved, its largest value among the retained
points, so a smaller $\rho$ admits only cleaner, higher-$k$ data. Enforcing this
data-dependent lower bound lets both filling regimes converge to the identical
target floor, and tightens the recovery band to $0.94\text{--}1.01\times$.

Two qualifications are load-bearing. First, the convergence is conditional on the
comparison budget actually reaching $k_{\min}$; a grid that cannot is flagged and
\emph{not fitted}, since a fit performed below its own window silently misreports
the floor rather than failing. That proviso is not free --- the empty filling's
larger $c_{\text{oracle}}$ raises $k_{\min}$ by nearly the same order of
magnitude. Second, what this eliminates is the need for a \emph{calibrated
constant per filling}, not the choice of filling itself. The closure decision
remains a modelling choice that determines $b_1$ and the dimension of the
harmonic space, and hence what the certificate is measuring at all: one
circulating flow on a single triangle reads as pure harmonic when the triple is
left open and pure curl when it is closed. The filling must still be declared;
what no longer has to be re-derived per filling is the estimation window.

\subsection{Quantization Artifacts and Quantized Emitters}
When mapping real-valued log-odds models to discrete comparison logs, two severe
edge cases surface at the quantization boundary. Both are dangerous for the same
reason: each yields a well-formed decomposition rather than an exception, so
nothing downstream can detect that the returned quantity is meaningless.

\paragraph{The Low-$k$ Tie Collapse Trap:} At the lower comparison limit
($k = 2$), passing a binary rule through a quantized emitter computes
$\text{round}(2 \cdot \sigma(1)) = 1$, generating an identical 1--1 tie on every
edge. This forces the entire emitted flow $Y \to 0$, rendering structural
asymmetries unobservable without triggering an explicit runtime exception.

\paragraph{Mixed-Configuration Rescaling Distortion:} Approximating a binary rule
by emitting all $R$ rows unidirectionally yields an edge magnitude of
$\pm\log(2R - 1)$. While exact for pure sign-rule configurations --- mass
fractions being scale-invariant --- executing this path within a mixed population
introduces a scaling mismatch that magnifies the complex block against adjacent
blocks by that same $\log(2R-1)$ factor: $2.708\times$ at the $R = 8$ at which
the effect was first measured, and $4.84\times$ at the instrument's present
default of $R = 64$. The resulting round-trip error of $0.269$ masquerades as an
intrinsic decomposition artifact.

The instrument suppresses both by dispatching per block across three
non-interchangeable emission paths rather than collapsing to a single one.
Native win-count replay is bit-exact, reproducing the internal decomposition to a
measured round-trip deviation of $0.00\mathrm{e}{+}00$. The unidirectional sign
path is retained but \emph{gated} to configurations that are entirely sign rules,
where a uniform rescaling cannot distort a mix. The quantized magnitude path
$w = \text{round}(k \cdot \sigma(Y))$ is approximate by construction --- exact
only in the limit $k \to \infty$ --- and reports its residual rather than
absorbing it into a tolerance. A flow that quantizes away entirely raises, while
an edge that legitimately rounds to a tie is counted and reported.
